\section{Multi-Agent Construction Architecture}

The main practical question is not whether such repositories are conceptually appealing, but whether they can be built. I argue that they can, using a divide-and-conquer multi-agent architecture.

Archaeologist agents would work over pre-digital and historical corpora: museum catalogs, excavation reports, archival documents, trade manuals, and anthropological studies [1, 2, 3, 4, 6, 8, 40]. Their role would be to recover early lineage structure and historically attested pressures.

Patent agents would extract formal lineage relations and inventor-stated problem definitions from patent corpora [15, 17, 19, 20, 37, 39]. These are especially valuable because many patents explicitly articulate the deficiency of prior designs.

Review-mining agents would process modern consumer and practitioner discourse: reviews, repair forums, technical journalism, complaint databases, and product communities [7, 9, 27, 28, 29, 36]. Their role would be to identify repeated user-expressed pressures and normalize them into structured annotations.

Diff agents would reconcile heterogeneous evidence into commit-level summaries. They would distinguish direct evidence from inference, manage conflicting accounts, and generate standardized change narratives.

Observer agents would operate across repositories to detect macro-patterns, especially convergence among independent lineages approaching similar solutions. These agents would be used retrospectively first and, if validated, prospectively for signal detection.

The architecture is intentionally modular. Different domains may require different source mixes, confidence policies, and extraction heuristics. The key point is not that one monolithic agent can solve the problem, but that heterogeneous, role-specific agents can collaboratively build lineage records that no single workflow could plausibly assemble.

To assess the feasibility of this architecture in a constrained setting, a small-scale pilot reconstruction was conducted on a historically documented artifact family.
